\chapter{Trừu Tượng Hoá Dữ Liệu}
\hypertarget{localtoc\thechapter}{}
\localtoc

\begin{problem}
    Hãy khai báo kiểu dữ liệu biểu diễn khái niệm phân số trong toán học và định nghĩa hàm nhập, hàm xuất cho kiểu dữ liệu này
\end{problem}
\begin{problem}
    Hãy khai báo kiểu dữ liệu biểu diễn khái niệm hỗn số trong toán học và định nghĩa hàm nhập, hàm xuất cho kiểu dữ liệu này
\end{problem}
\begin{problem}
    Hãy khai báo kiểu dữ liệu biểu diễn khái niệm điểm trong mặt phẳng Oxy và định nghĩa hàm nhập, hàm xuất cho kiểu dữ liệu này
\end{problem}
\begin{problem}
    Hãy khai báo kiểu dữ liệu biểu diễn khái niệm điểm trong không gian Oxyz và định nghĩa hàm nhập, hàm xuất cho kiểu dữ liệu này
\end{problem}
\begin{problem}
    Hãy khai báo kiểu dữ liệu biểu diễn khái niệm đơn thức $P(x) = ax^n$ trong toán học và định nghĩa và định nghĩa hàm nhập, hàm xuất cho kiểu dữ liệu này
\end{problem}
\begin{problem}
    Hãy khai báo kiểu dữ liệu để biểu diễn khái niệm đa thức 1 biến trong toán học: $P(x) = a.n.X^n + a.(n-1).X^{n-1} + … + a.1.X + a.0$ và định nghĩa hàm nhập và hàm xuất cho kiểu dữ liệu này
\end{problem}
\begin{problem}
    Hãy khai báo kiểu dữ liệu biểu diễn ngày trong thế giới thực và định nghĩa hàm nhập, hàm xuất cho kiểu dữ liệu này
\end{problem}
\begin{problem}
    Hãy khai báo kiểu dữ liệu biểu diễn khái niệm đường thẳng ax + by + c = 0 trong mặt phẳng Oxy và định nghĩa hàm nhập, hàm xuất cho kiểu dữ liệu này
\end{problem}
\begin{problem}
    Hãy khai báo kiểu dữ liệu biểu diễn khái niệm đường tròn trong mặt phẳng Oxy và định nghĩa hàm nhập, hàm xuất cho kiểu dữ liệu này
\end{problem}
\begin{problem}
    Viết chương trình nhập tọa độ tâm và bán kính của 1 đường tròn trong mặt phẳng Oxy. Tính diện tích và chu vi của nó và xuất ra kết quả
\end{problem}
\begin{problem}
     Viết chương trình nhập tọa độ 3 đỉnh của 1 tam giác trong mặt phẳng Oxy. Tính diện tích, chu vi và tọa độ trọng tâm của tam giác và xuất ra kết quả
\end{problem}

\section{Đơn thức}

\begin{problem}
    Tính tích 2 đơn thức
\end{problem}
\begin{problem}
    Tính đạo hàm cấp 1 đơn thức
\end{problem}
\begin{problem}
    Tính thương 2 đơn thức
\end{problem}
\begin{problem}
    Tính đạo hàm cấp k đơn thức
\end{problem}
\begin{problem}
    Tính giá trị đơn thức tại vị trí $x=x0$
\end{problem}
\begin{problem}
    Định nghĩa toán tử (operator $*$) cho 2 đơn thức
\end{problem}
\begin{problem}
    Định nghĩa toán tử (operator$ /$) cho 2 đơn thức
\end{problem}

\section{Đa thức}

\begin{problem}
    Tính hiệu 2 đa thức
\end{problem}
\begin{problem}
    Tính tổng 2 đa thức
\end{problem}
\begin{problem}
    Tính tích 2 đa thức
\end{problem}
\begin{problem}
    Tính thương 2 đa thức
\end{problem}
\begin{problem}
    Tính đa thức dư của phép chia đa thức thứ nhất cho đa thức thứ hai
\end{problem}
\begin{problem}
    Tính đạo hàm cấp 1 của đa thức
\end{problem}
\begin{problem}
    Tính đạo hàm cấp k của đa thức
\end{problem}
\begin{problem}
    Tính giá trị của đa thức tại vị trí $x = x0$
\end{problem}
\begin{problem}
    Định nghĩa toán tử cộng (operator $+$) cho hai đa thức
\end{problem}
\begin{problem}
    Định nghĩa toán tử trừ (operator $-$) cho hai đa thức
\end{problem}
\begin{problem}
    Định nghĩa toán tử nhân (operator $*$) cho hai đa thức
\end{problem}
\begin{problem}
    Định nghĩa toán tử thương (operator $/$) cho hai đa thức
\end{problem}
\begin{problem}
    Tìm nghiệm của đa thức trong đoạn $[a, b]$ cho trước
\end{problem}

\section{Phân số}

\begin{problem}
    Rút gọn phân số
\end{problem}
\begin{problem}
    Tính tổng 2 phân số
\end{problem}
\begin{problem}
    Tính hiệu 2 phân số
\end{problem}
\begin{problem}
    Tính tích 2 phân số
\end{problem}
\begin{problem}
    Tính thương 2 phân số
\end{problem}
\begin{problem}
    Kiểm tra phân số tối giản
\end{problem}
\begin{problem}
    Qui đồng phân số
\end{problem}
\begin{problem}
    Kiểm tra phân số dương
\end{problem}
\begin{problem}
    Kiểm tra phân số âm
\end{problem}
\begin{problem}
    So sánh 2 phân số: hàm trả về 1 trong 3 giá trị: 0,-1,1
\end{problem}
\begin{problem}
    Định nghĩa toán tử operator $+$ cho 2 phân số
\end{problem}
\begin{problem}
    Định nghĩa toán tử operator $-$ cho 2 phân số
\end{problem}
\begin{problem}
    Định nghĩa toán tử operator $*$ cho 2 phân số
\end{problem}
\begin{problem}
    Định nghĩa toán tử operator $/$ cho 2 phân số
\end{problem}
\begin{problem}
    Định nghĩa toán tử operator $++$ cho 2 phân số
\end{problem}
\begin{problem}
    Định nghĩa toán tử operator $—$ cho 2 phân số
\end{problem}
\begin{problem}
    
\end{problem}
\begin{problem}
    
\end{problem}

\section{Hỗn số}

\begin{problem}
    Khai báo dữ liệu để biểu diễn thông tin của một hỗn số
\end{problem}
\begin{problem}
    Nhập hỗn số
\end{problem}
\begin{problem}
    Xuất hỗn số
\end{problem}
\begin{problem}
    Rút gọn hỗn số
\end{problem}
\begin{problem}
    Tính tổng 2 hỗn số
\end{problem}
\begin{problem}
    Tính hiệu 2 hỗn số
\end{problem}
\begin{problem}
    Tính tích 2 hỗn số
\end{problem}
\begin{problem}
    Tính thương 2 hỗn số
\end{problem}
\begin{problem}
    Kiểm tra hỗn số tối giản
\end{problem}
\begin{problem}
    Qui đồng 2 hỗn số
\end{problem}


\section{Số phức}

\begin{problem}
    Khai báo biểu diễn thông tin số phức
\end{problem}
\begin{problem}
    Nhập số phức
\end{problem}
\begin{problem}
    Xuất số phức
\end{problem}
\begin{problem}
    Tính tổng 2 số phức
\end{problem}
\begin{problem}
    Tính hiệu 2 số phức
\end{problem}
\begin{problem}
    Tính tích 2 số phức
\end{problem}
\begin{problem}
    Tính thương 2 số phức
\end{problem}
\begin{problem}
    Tính lũy thừa bậc n số phức
\end{problem}

\section{Điểm trong mặt phẳng Oxy}

\begin{problem}
    khai báo dữ liệu điểm OXY
\end{problem}
\begin{problem}
    Nhập tọa độ điểm trong mặt phẳng
\end{problem}
\begin{problem}
    Xuất tọa độ điểm trong mặt phẳng
\end{problem}
\begin{problem}
    Tính khoảng cách giữa 2 điểm
\end{problem}
\begin{problem}
    Tính khoảng cách 2 điểm theo phương Ox
\end{problem}
\begin{problem}
    Tính khoảng cách 2 điểm theo phương Oy
\end{problem}
\begin{problem}
    Tìm tọa độ điểm đối xứng qua gốc tọa độ
\end{problem}
\begin{problem}
    Tìm điểm đối xứng qua trục hoành
\end{problem}
\begin{problem}
    Tìm điểm đối xứng qua trục tung
\end{problem}
\begin{problem}
    Tìm điểm đối xứng qua đường phân giác thứ 1 $(y=x)$
\end{problem}
\begin{problem}
    Tìm điểm đối xứng qua đường phân giác thứ 2 $(y=-x)$
\end{problem}
\begin{problem}
    Kiểm tra điểm thuộc phần tư thứ 1 không?
\end{problem}
\begin{problem}
    Kiểm tra điểm thuộc phần tư thứ 2 không?
\end{problem}
\begin{problem}
    Kiểm tra điểm thuộc phần tư thứ 3 không?
\end{problem}
\begin{problem}
    Kiểm tra điểm thuộc phần tư thứ 4 không?
\end{problem}

\section{Điểm trong không gian Oxyz}

\begin{problem}
    Khai báo kiểu dữ liệu biểu diễn tọa độ điểm trong không gian Oxyz
\end{problem}
\begin{problem}
    Nhập tọa độ điểm trong không gian Oxyz
\end{problem}
\begin{problem}
    Xuất tọa độ điểm theo định dạng $(x, y, z)$
\end{problem}
\begin{problem}
    Tính khoảng cách giữa 2 điểm trong không gian
\end{problem}
\begin{problem}
    Tính khoảng cách giữa 2 điểm trong không gian theo phương x
\end{problem}
\begin{problem}
    Tính khoảng cách giữa 2 điểm trong không gian theo phương y
\end{problem}
\begin{problem}
    Tính khoảng cách giữa 2 điểm trong không gian theo phương z
\end{problem}
\begin{problem}
    Tìm tọa độ điểm đối xứng qua gốc tọa độ
\end{problem}
\begin{problem}
    Tìm tọa độ điểm đối xứng qua mặt phẳng Oxy
\end{problem}
\begin{problem}
    Tìm tọa độ điểm đối xứng qua mặt phẳng Oxz
\end{problem}
\begin{problem}
    Tìm tọa độ điểm đối xứng qua mặt phẳng Oyz
\end{problem}

\section{Đường tròn trong mặt phẳng Oxy}

\begin{problem}
    Khai báo kiểu dữ liệu để biểu diễn đường tròn
\end{problem}
\begin{problem}
    Nhập đường tròn
\end{problem}
\begin{problem}
    Xuất đường tròn theo định dạng $((x, y), r)$
\end{problem}
\begin{problem}
    Tính chu vi đường tròn
\end{problem}
\begin{problem}
    Tính diện tích đường tròn
\end{problem}
\begin{problem}
    Xét vị trí tương đối giữa 2 đường tròn( không cắt nhau, tiếp xúc, cắt nhau)
\end{problem}
\begin{problem}
    Kiểm tra 1 tọa độ điểm có nằm trong đường tròn hay không
\end{problem}
\begin{problem}
    Cho 2 đường tròn. Tính diện tích phần mặt phẳng bị phủ bởi 2 đường tròn đó
\end{problem}

\section{Hình cầu trong không gian Oxyz}

\begin{problem}
    Khai báo kiểu dữ liệu để biểu diễn hình cầu trong không gian Oxyz
\end{problem}
\begin{problem}
    Nhập hình cầu
\end{problem}
\begin{problem}
    Xuất hình cầu theo định dạng $((x, y, z), r)$
\end{problem}
\begin{problem}
    Tính diện tích xung quanh hình cầu
\end{problem}
\begin{problem}
    Tính thể tích hình cầu
\end{problem}
\begin{problem}
    Xét vị trí tương đối giữa 2 hình cầu(không cắt nhau, tiếp xúc, cắt nhau)
\end{problem}
\begin{problem}
    Kiểm tra 1 tọa độ điểm có nằm bên trong hình cầu hay không
\end{problem}


\section{Tam giác trong mặt phẳng Oxy}

\begin{problem}
    Khai báo kiểu dữ liệu để biểu diễn tam giác trong mặt phẳng Oxy
\end{problem}
\begin{problem}
    Nhập tam giác
\end{problem}
\begin{problem}
    Xuất tam giác theo định dạng $((x1, y1); (x2, y2); (x3, y3))$
\end{problem}
\begin{problem}
    Kiểm tra tọa độ 3 đỉnh có thật sự lập thành 3 đỉnh của 1 tam giác không
\end{problem}
\begin{problem}
    Tính chu vi tam giác
\end{problem}
\begin{problem}
    Tính diện tích tam giác
\end{problem}
\begin{problem}
    Tìm tọa độ trọng tâm tam giác
\end{problem}
\begin{problem}
    Tìm 1 đỉnh trong tam giác có hoành độ lớn nhất
\end{problem}
\begin{problem}
    Tìm 1 đỉnh trong tam giác có tung độ nhỏ nhất
\end{problem}
\begin{problem}
    Tính tổng khoảng cách từ điểm $P(x, y)$ tới 3 đỉnh của tam giác
\end{problem}
\begin{problem}
    Kiểm tra 1 tọa độ điểm có nằm trong tam giác hay không
\end{problem}
\begin{problem}
    Hãy cho biết dạng của tam giác(đều, vuông, vuông cân, cân, thường)
\end{problem}

\section{Ngày}

\begin{problem}
    Khai báo kiểu dữ liệu để biểu diễn ngày
\end{problem}
\begin{problem}
    Nhập ngày
\end{problem}
\begin{problem}
    Xuất ngày theo định dạng (ng/th/nm)
\end{problem}
\begin{problem}
    Kiểm tra năm nhuận
\end{problem}
\begin{problem}
    Tính số thứ tự ngày trong năm
\end{problem}
\begin{problem}
    Tính số thứ tự ngày kể từ ngày 1/1/1
\end{problem}
\begin{problem}
    Tìm ngày khi biết năm và số thứ tự của ngày trong năm
\end{problem}
\begin{problem}
    Tìm ngày khi biết số thứ tự ngày kể từ ngày 1/1/1
\end{problem}
\begin{problem}
    Tìm ngày kế tiếp
\end{problem}
\begin{problem}
    Tìm ngày hôm qua
\end{problem}
\begin{problem}
    Tìm ngày kế đó k ngày
\end{problem}
\begin{problem}
    Tìm ngày trước đó k ngày
\end{problem}
\begin{problem}
    Khoảng cách giữa 2 ngày
\end{problem}
\begin{problem}
    So sánh 2 ngày

    Tính thứ của ngày bất kỳ trong năm(Dùng CT Zeller)
\end{problem}
\begin{problem}
    Hãy khai báo kiểu dữ liệu để biểu diễn thông tin của 1 tỉnh (TINH). Biết rằng một tỉnh gồm những thành phần sau:

Mã tỉnh: Kiểu số nguyên 2 byte

Tên tỉnh: Chuỗi tối đa 30 ký tự

Diện tích: Kiểu số thực

Sau đó viết hàm nhập, xuất cho kiểu dữ liệu này
\end{problem}
\begin{problem}
    Hãy khai báo kiểu dữ liệu để biểu diễn thông tin của một hộp sữa (HOPSUA). Biết rằng một hộp sữa gồm các thành phần sau:

Nhãn hiệu: chuỗi tối đa 20 ký tự

Trọng lượng: kiểu số thực

Hạn sử dụng: Kiểu dữ liệu Ngày (NGAY)

Sau đó viết hàm nhập, xuất cho kiểu dữ liệu này
\end{problem}
\begin{problem}
    Hãy khai báo kiểu dữ liệu để biểu diễn thông tin của 1 vé xem phim (VE). Biết rằng 1 vé xem phim gồm những thành phần sau:

Tên phim: Chuỗi tối đa 20 ký tự

Giá tiền: kiểu số nguyên 4 byte

Xuất chiếu: kiểu thời gian (THOIGIAN)

Ngày xem: kiểu dữ liệu ngày (NGAY)

Sau đó viết hàm nhập, xuất cho kiểu dữ liệu này
\end{problem}
\begin{problem}
    Hãy khai báo kiểu dữ liệu để biểu diễn thông tin của một mặt hang (MATHANG). Biết rằng một mặt hang gồm những thành phần sau:

Tên mặt hàng: chuỗi tối đa 20 ký tự

Đơn giá: kiểu số nguyên 4 byte

Số lượng tồn: kiểu số nguyên 4 byte

Sau đó viết hàm nhập, xuất cho kiểu dữ liệu này

\end{problem}
\begin{problem}
    
Bài 611: Hãy khai báo kiểu dữ liệu để biểu diễn thông tin của một chuyến bay. Biết rằng một chuyến bay gồm những thành phần sau:

Mã chuyến bay: chuỗi tối đa 5 ký tự

Ngày bay: kiểu dữ liệu ngày

Giờ bay: kiểu thời gian

Nơi đi: chuỗi tối đa 20 ký tự

Nơi đến: chuỗi tối đa 20 ký tự

Sau đó viết hàm nhập, xuất cho kiểu dữ liệu này
\end{problem}
\begin{problem}
    Hãy khai báo kiểu dữ liệu để biểu diễn thông tin của một cầu thủ. Biết rằng một cầu thủ gồm những thành phần sau:

Mã cầu thủ: chuỗi tối đa 10 ký tự

Tên cầu thủ: chuỗi tối đa 30 ký tự

Ngày sinh: kiểu dữ liệu ngày

Sau đó viết hàm nhập, xuất cho kiểu dữ liệu này
\end{problem}
\begin{problem}
    Hãy khai báo kiểu dữ liệu để biểu diễn thông tin của một đội bóng. (DOIBONG). Biết rằng một đội bóng gồm những thành phần sau:

Mã đội bóng: chuỗi tối đa 5 ký tự

Tên đội bóng: chuỗi tối đa 30 ký tự

Danh sách các cầu thủ: mảng 1 chiều các cầu thủ (tối đa 30 phần tử)

Sau đó viết hàm nhập, xuất cho kiểu dữ liệu này
\end{problem}
\begin{problem}
    Hãy khai báo kiểu dữ liệu để biểu diễn thông tin của một nhân viên (NHANVIEN). Biết rằng một nhân viên gồm những thành phần sau:

Mã nhân viên: chuỗi tối đa 5 ký tự

Tên nhân viên: chuỗi tối đa 30 ký tự

Lương nhân viên: kiểu số thực

Sau đó viết hàm nhập, xuất cho kiểu dữ liệu này

\end{problem}
\begin{problem}
    Hãy khai báo kiểu dữ liệu để biểu diễn thông tin của một thí sinh (THISINH). Biết rằng một thí sinh gồm những thành phần sau:

Mã thí sinh: chuỗi tối đa 5 ký tự

Họ tên thí sinh: chuỗi tối đa 30 ký tự

Điểm toán: kiểu số thực

Điểm lý: kiểu số thực

Điểm hóa: kiểu số thực

Điểm tổng cộng: kiểu số thực

Sau đó viết hàm nhập, xuất cho kiểu dữ liệu này
\end{problem}
\begin{problem}
    Hãy khai báo kiểu dữ liệu để biểu diễn thông tin của một luận văn (LUANVAN). Biết rằng một luận văn gồm những thành phần sau:

Mã luận văn: chuỗi tối đa 10 ký tự

Tên luận văn: chuỗi tối đa 100 ký tự

Họ tên sinh viên thực hiện: chuỗi tối đa 30 ký tự

Họ tên giảng viên hướng dẫn: chuỗi tối đa 30 ký tự

Năm thực hiện: kiểu số nguyên 2 byte

Sau đó viết hàm nhập, xuất cho kiểu dữ liệu này
\end{problem}
\begin{problem}
    Hãy khai báo kiểu dữ liệu để biểu diễn thông tin của một học sinh (HOCSINH). Biết rằng một lớp học gồm những thành phần sau:

Tên học sinh: chuỗi tối đa 30 ký tự

Điểm toán: kiểu số nguyên 2 byte

Điểm văn: kiểu số nguyên 2 byte

Điểm trung bình: kiểu số thực

Sau đó viết hàm nhập, xuất cho kiểu dữ liệu này
\end{problem}
\begin{problem}
    Hãy khai báo kiểu dữ liệu để biểu diễn thông tin của một lớp học (LOPHOC). Biết rằng một lớp học gồm những thành phần sau:

Tên lớp: chuỗi tối đa 30 ký tự

Sĩ số: kiểu số nguyên 2 byte

Danh sách các học sinh trong lớp ( tối đa 50 học sinh)

Sau đó viết hàm nhập, xuất cho kiểu dữ liệu này
\end{problem}
\begin{problem}
    Hãy khai báo kiểu dữ liệu để biểu diễn thông tin của một sổ tiết kiệm (SOTIETKIEM). Biết rằng một sổ tiết kiệm gồm những thành phần sau:

Mã sổ: chuỗi tối đa 5 ký tự

Loại tiết kiệm: chuỗi tối đa 10 ký tự

Họ tên khách hàng: chuỗi tối đa 30 ký tự

Chứng minh nhân dân: kiểu số nguyên 4 byte

Ngày mở sổ: kiểu dữ liệu ngày

Số tiền gửi: kiểu số thực

Sau đó viết hàm nhập, xuất cho kiểu dữ liệu này
\end{problem}
\begin{problem}
    Hãy khai báo kiểu dữ liệu để biểu diễn thông tin của một đại lý (DAILY). Biết rằng một đại lý gồm những thành phần sau:

Mã đại lý: chuỗi tối đa 5 ký tự

Tên đại lý: chuỗi tối đa 30 ký tự

Điện thoại: kiểu số nguyên 4 byte

Ngày tiếp nhận: kiểu dữ liệu ngày

Địa chỉ: chuỗi tối đa 50 ký tự

E-Mail: chuỗi tối đa 50 ký tự

Sau đó viết hàm nhập, xuất cho kiểu dữ liệu này
\end{problem}

